%=============================================================
% Capitolo 3 — Metodi Diretti per Sistemi Lineari
%=============================================================
\chapter{Metodi Diretti per Sistemi Lineari}
\label{ch:diretti}

\section{Panoramica}

I \textbf{metodi diretti} risolvono il sistema $Ax = b$ in un numero finito di
operazioni aritmetiche, producendo (in aritmetica esatta) la soluzione esatta.
Il costo è tipicamente $\mathcal{O}(n^3)$ operazioni in virgola mobile.

L'idea centrale è \textbf{fattorizzare} la matrice $A$ nel prodotto di matrici più
semplici — tipicamente triangolari o diagonali — per ridursi a sistemi facilmente
risolvibili.

%------------------------------------------------------------
\section{Sistemi triangolari}
\label{sec:triangolari}
%------------------------------------------------------------

\subsection{Sistema triangolare inferiore}

Un sistema $Lx = b$ con $L$ triangolare inferiore ($l_{ij} = 0$ per $j > i$) si
risolve con la \textbf{sostituzione in avanti} (\emph{forward substitution}):
\begin{equation}
  x_i = \frac{1}{l_{ii}}\left(b_i - \sum_{j=1}^{i-1} l_{ij}\, x_j\right),
  \quad i = 1, 2, \ldots, n.
  \label{eq:fwd_sub}
\end{equation}

\subsection{Sistema triangolare superiore}

Un sistema $Ux = b$ con $U$ triangolare superiore ($u_{ij} = 0$ per $j < i$) si
risolve con la \textbf{sostituzione all'indietro} (\emph{back substitution}):
\begin{equation}
  x_i = \frac{1}{u_{ii}}\left(b_i - \sum_{j=i+1}^{n} u_{ij}\, x_j\right),
  \quad i = n, n-1, \ldots, 1.
  \label{eq:bck_sub}
\end{equation}

Entrambe le sostituzioni hanno costo $\mathcal{O}(n^2)$.

%------------------------------------------------------------
\section{Fattorizzazione LU}
\label{sec:lu}
%------------------------------------------------------------

\begin{definition}[Fattorizzazione LU]
Una \emph{fattorizzazione LU} di $A \in \R^{n \times n}$ è la scrittura
\begin{equation}
  A = LU,
  \label{eq:LU}
\end{equation}
dove $L$ è triangolare inferiore con elementi diagonali unitari ($l_{ii} = 1$) e
$U$ è triangolare superiore.
\end{definition}

Con la fattorizzazione LU il sistema $Ax = b$ si risolve in due passi:
\begin{align}
  Ly &= b \quad \text{(forward substitution)}, \label{eq:forw}\\
  Ux &= y \quad \text{(back substitution)}.  \label{eq:back}
\end{align}
Il costo totale è $\mathcal{O}(n^3)$ per la fattorizzazione e $\mathcal{O}(n^2)$ per
ciascuna risoluzione successiva (utile quando si deve risolvere più sistemi con la
stessa matrice $A$).

%------------------------------------------------------------
\section{Eliminazione di Gauss}
\label{sec:gauss}
%------------------------------------------------------------

L'\textbf{eliminazione di Gauss} è l'algoritmo classico per calcolare la
fattorizzazione LU. A ogni passo $k$ si elimina l'incognita $x_k$ dalle righe
$k+1, \ldots, n$ mediante opportune combinazioni lineari.

\subsection{Un passo dell'eliminazione}

Al passo $k$ si ha a disposizione la matrice parzialmente eliminata
$A^{(k)} = (a^{(k)}_{ij})$. Per ogni riga $i > k$ si calcola il
\emph{moltiplicatore}:
\begin{equation}
  m_{ik} = \frac{a^{(k)}_{ik}}{a^{(k)}_{kk}},
  \label{eq:moltiplicatore}
\end{equation}
e si aggiorna:
\begin{equation}
  a^{(k+1)}_{ij} = a^{(k)}_{ij} - m_{ik}\, a^{(k)}_{kj},
  \quad j = k, k+1, \ldots, n.
  \label{eq:update_gauss}
\end{equation}

L'elemento $a^{(k)}_{kk}$ si chiama \textbf{pivot} al passo $k$.

\begin{algorithm}[htbp]
  \caption{Eliminazione di Gauss (fattorizzazione LU)}
  \label{alg:gauss}
  \begin{algorithmic}[1]
    \Require $A \in \R^{n \times n}$ non singolare
    \Ensure $L, U$ tali che $A = LU$
    \State $U \gets A$; \; $L \gets I_n$
    \For{$k = 1$ \textbf{to} $n-1$}
      \For{$i = k+1$ \textbf{to} $n$}
        \State $m_{ik} \gets u_{ik} / u_{kk}$
        \State $l_{ik} \gets m_{ik}$
        \For{$j = k$ \textbf{to} $n$}
          \State $u_{ij} \gets u_{ij} - m_{ik} \cdot u_{kj}$
        \EndFor
      \EndFor
    \EndFor
    \State \Return $L$, $U$
  \end{algorithmic}
\end{algorithm}

Il costo dell'eliminazione di Gauss è $\frac{2}{3}n^3 + \mathcal{O}(n^2)$
operazioni in virgola mobile.

\begin{theorem}[Esistenza della fattorizzazione LU]
La fattorizzazione LU senza pivoting esiste ed è unica se e solo se tutti i
\emph{minori principali} di $A$ sono non nulli.
\end{theorem}

%------------------------------------------------------------
\section{Pivoting}
\label{sec:pivoting}
%------------------------------------------------------------

Nella pratica il pivot $a^{(k)}_{kk}$ può essere zero (impossibilità) o molto
piccolo (instabilità numerica). Per evitare questi problemi si adotta il
\textbf{pivoting parziale per colonne}: ad ogni passo $k$ si sceglie come pivot
l'elemento di modulo massimo nella colonna $k$ tra le righe $k, k+1, \ldots, n$,
e si scambiano le righe corrispondentemente.

Con il pivoting parziale si ha la fattorizzazione:
\begin{equation}
  PA = LU,
  \label{eq:PLU}
\end{equation}
dove $P$ è una \textbf{matrice di permutazione} che registra gli scambi di riga.

\begin{remark}[Stabilità]
L'eliminazione di Gauss con pivoting parziale è numericamente stabile per quasi
tutte le matrici in pratica, sebbene esistano esempi patologici di instabilità.
Il pivoting totale (su righe e colonne) è più stabile ma meno usato per il
maggiore costo computazionale.
\end{remark}

%------------------------------------------------------------
\section{Matrici sparse e metodi a basso costo}
\label{sec:sparse}
%------------------------------------------------------------

Quando $A$ ha molti elementi nulli (\emph{matrice sparsa}), la fattorizzazione LU
standard non sfrutta la struttura e può produrre un \emph{fill-in} (riempimento
degli zeri) costoso. Tecniche specializzate (reordering, fattorizzazioni incomplete)
permettono di gestire efficientemente sistemi sparsi di grandi dimensioni.

%------------------------------------------------------------
\section{Costo computazionale e confronto}
\label{sec:costo_diretti}
%------------------------------------------------------------

\begin{table}[htbp]
  \centering
  \caption{Costo computazionale dei principali metodi diretti.}
  \label{tab:costo}
  \begin{tabular}{lcc}
    \toprule
    \textbf{Operazione} & \textbf{Flop (ordine)} & \textbf{Note} \\
    \midrule
    Fattorizzazione LU             & $\tfrac{2}{3}n^3$ & una sola volta per $A$ \\
    Forward/back substitution      & $n^2$             & per ogni $b$ \\
    Calcolo esplicito di $A^{-1}$  & $2n^3$            & sconsigliato \\
    \bottomrule
  \end{tabular}
\end{table}
